%% Question bank chapter file - generated by bank_tools.py build (do not edit while the bank is still being built:
%% it is regenerated from the per-test data in Question Bank\data). Every question is one bankq block:
%%   \begin{bankq}[id=..., source={...}, paper=.., qno=.., type=..., subtopic={...}, answer={...}, answer_src=key|solved,
%%                 working={...}, difficulty=1-10, diff_parts={K S I C F}, diff_note={...},
%%                 group=<passage id or empty>, also_in={duplicate copies in other tests}]
%% difficulty = min(10, K+S+I+C+F-1): K concepts 1-3, S steps 1-4, I insight 0-2, C computation 0-2, F format 0-1
%%   ...question LaTeX, options, figures (\bfig / \bfigi -> img/)...
%%   \end{bankq}
%% Passages shared by several questions are bankpassage blocks; their questions carry group=<passage id>.
\documentclass[10pt]{article}
\usepackage[a4paper,margin=1.8cm]{geometry}
\usepackage[T1]{fontenc}
\usepackage{amsmath,amssymb,mathtools}
\usepackage{graphicx}
\usepackage[export]{adjustbox}
\usepackage{xcolor,needspace,keyval}
\usepackage{enumitem}
\usepackage{booktabs,array,tabularx}
\usepackage[hidelinks]{hyperref}
\setlength{\parindent}{0pt}
\setlength{\parskip}{3pt}
\makeatletter
\define@key{bq}{id}{\def\bq@id{#1}}\define@key{bq}{source}{\def\bq@src{#1}}\define@key{bq}{paper}{\def\bq@paper{#1}}
\define@key{bq}{qno}{\def\bq@qno{#1}}\define@key{bq}{type}{\def\bq@type{#1}}\define@key{bq}{subtopic}{\def\bq@sub{#1}}
\define@key{bq}{answer}{\def\bq@ans{#1}}\define@key{bq}{answer_src}{\def\bq@asrc{#1}}\define@key{bq}{working}{\def\bq@work{#1}}
\define@key{bq}{group}{\def\bq@group{#1}}\define@key{bq}{also_in}{\def\bq@also{#1}}
\define@key{bq}{difficulty}{\def\bq@diff{#1}}\define@key{bq}{diff_parts}{\def\bq@dparts{#1}}\define@key{bq}{diff_note}{\def\bq@dnote{#1}}
\newenvironment{bankq}[1][]{\def\bq@work{}\def\bq@also{}\def\bq@group{}\def\bq@diff{?}\def\bq@dparts{}\def\bq@dnote{}\setkeys{bq}{#1}%
  \par\needspace{4\baselineskip}\medskip\noindent{\small\textbf{\bq@id}\hfill\color{gray}[\bq@src, Paper \bq@paper, Q\bq@qno]
  \bq@type\ $\cdot$ \bq@sub\ $\cdot$ \textbf{D\bq@diff}}\par\nopagebreak\relax}
  {\par{\small\color{teal}Answer: \textbf{\bq@ans}\ifx\bq@work\@empty\else\ (solved: \bq@work)\fi
  \ifx\bq@also\@empty\else\par\color{gray}Also in: \bq@also\fi}\par}
\define@key{bp}{id}{\def\bp@id{#1}}\define@key{bp}{source}{\def\bp@src{#1}}\define@key{bp}{questions}{\def\bp@qs{#1}}
\newenvironment{bankpassage}[1][]{\setkeys{bp}{#1}\par\needspace{4\baselineskip}\bigskip\noindent
  {\small\textbf{Passage \bp@id}\hfill\color{gray}[\bp@src; questions \bp@qs]}\par\nopagebreak\relax\itshape}{\par}
\makeatother
\newcommand{\bfig}[2]{\par\smallskip{\centering\includegraphics[width=#1pt,max width=0.62\linewidth,max height=0.32\textheight]{img/#2}\par}\smallskip}
\newcommand{\bfigi}[2]{\includegraphics[width=#1pt,max width=\linewidth,max height=0.32\textheight,valign=c]{img/#2}}
\newcommand{\bmissing}[1]{\textit{[figure missing in the source conversion: \detokenize{#1}]}}
\newenvironment{qoptions}{\begin{itemize}[leftmargin=1.8em,itemsep=1pt,topsep=2pt,label={}]}{\end{itemize}}
\newcommand{\opt}[2]{\item[(#1)] #2}
\providecommand{\fwstart}{}\providecommand{\fwend}{}\providecommand{\qfullsplit}{}
\newcolumntype{L}[1]{>{\raggedright\arraybackslash}p{#1}}
\newcolumntype{C}[1]{>{\centering\arraybackslash}p{#1}}
\newcolumntype{R}[1]{>{\raggedleft\arraybackslash}p{#1}}
\providecommand{\sectionhdr}[1]{}\providecommand{\subsectionhdr}[1]{}


\begin{document}
{\Large\bfseries Monotonicity – Max/Min --- Question Bank}\par\smallskip
{\small Mathematics chapter M25. Questions from all converted test papers, grouped by source series. Answers: printed key where it matches the paper, otherwise solved (answer\_src=solved, with working). Exact repeats across papers appear once, with also\_in listing the other copies.}\par\medskip\hrule

\section*{Score-I papers}
\begin{bankq}[id=SI12-P1-MATH-S1-Q07, source={Score-I - Test 12 (09 Mar 2022)}, paper=1, qno=7, type=multi, subtopic={Extrema}, answer={A,C,D}, answer_src=key, difficulty=5, diff_parts={K1 S3 I1 C0 F1}, diff_note={Sign analysis of f' near x=1 and x=2 for even/odd n}]
Let $f(x)=(x-1)^4(x-2)^n$, $n\in N$. $f(x)$ has
\begin{qoptions}
\opt{A}{Local minimum at $x=2$ if $n$ is even}
\opt{B}{Local minimum at $x=1$ if $n$ is odd}
\opt{C}{Local maximum at $x=1$ if $n$ is odd}
\opt{D}{Local minimum at $x=1$ if $n$ is even}
\end{qoptions}
\end{bankq}
\begin{bankq}[id=SI12-P1-MATH-S2-Q01, source={Score-I - Test 12 (09 Mar 2022)}, paper=1, qno=1, type=numerical, subtopic={Increasing and decreasing functions}, answer={33.33}, answer_src=key, difficulty=7, diff_parts={K2 S3 I1 C1 F1}, diff_note={f'=4 sin^2x (cos x + a) via sum-to-product, need a>=1 so lambda=1}]
The values of `a' for which the function $f(x)=\sin x-a\sin2x-\dfrac13\sin3x+2ax$ increases throughout the number line is $[\lambda,\infty)$, then find $\dfrac{100\lambda}3$
\end{bankq}

\section*{TOAS papers}
\begin{bankpassage}[id=TOAS04-P1-MATH-S3-G1-2, source={TOAS - Test 04 (15 Mar 2022)}, questions={1-2}]
If $f(x)=4x^3-x^2-2x+1$ and \[g(x)=\left\{\begin{array}{ll}\min\{f(t);\ 0\leqslant t\leqslant x\} & ;\ \ 0\leqslant x\leqslant1\\[2pt] 3-x & ;\ \ 1<x\leqslant2\end{array}\right\},\] then
\end{bankpassage}
\begin{bankq}[id=TOAS04-P1-MATH-S3-Q01, source={TOAS - Test 04 (15 Mar 2022)}, paper=1, qno=1, type=numerical, subtopic={Extrema}, answer={0.50}, answer_src=key, difficulty=4, diff_parts={K1 S2 I1 C0 F1}, diff_note={f decreasing on [0,1/4] so g(1/4)=f(1/4)}, group=TOAS04-P1-MATH-S3-G1-2]
$g\left(\dfrac14\right)$ is equal to
\end{bankq}
\begin{bankq}[id=TOAS04-P1-MATH-S3-Q02, source={TOAS - Test 04 (15 Mar 2022)}, paper=1, qno=2, type=numerical, subtopic={Extrema}, answer={0.25}, answer_src=key, difficulty=5, diff_parts={K1 S3 I1 C0 F1}, diff_note={Interior minimum of f at x=1/2 gives g(3/4)=f(1/2)}, group=TOAS04-P1-MATH-S3-G1-2]
$g\left(\dfrac34\right)$ is equal to
\end{bankq}
\begin{bankq}[id=TOAS08-P1-MATH-S1-Q02, source={TOAS - Test 08 (15 May 2022)}, paper=1, qno=2, type=multi, subtopic={Extrema}, answer={C}, answer_src=key, difficulty=7, diff_parts={K2 S4 I1 C0 F1}, diff_note={Cubic from extremum conditions and values; greatest-integer of the roots}]
Let $f(x)$ be a cubic polynomial which has local maximum at $x=-1$ and $f'(x)$ has local minimum at $x=1$. If $f(-1)=10$ and $f(0)=5$ and roots of $f(x)$ are $\alpha,\beta,\gamma$, then $[\alpha]^2+[\beta]^2+[\gamma]^2$ is (where $[.]$ denotes greatest integer function)
\begin{qoptions}
\opt{A}{20}
\opt{B}{22}
\opt{C}{25}
\opt{D}{31}
\end{qoptions}
\end{bankq}
\begin{bankq}[id=TOAS08-P2-MATH-S1-Q03, source={TOAS - Test 08 (15 May 2022)}, paper=2, qno=3, type=multi, subtopic={Extrema}, answer={A,B,C,D}, answer_src=key, difficulty=7, diff_parts={K2 S4 I1 C0 F1}, diff_note={Cubic from extremum conditions: roots, intervals and range}]
Let $f(x)$ be a cubic polynomial which is having local maximum at $(1,2)$ and $f'(x)$ has local extremum at $x=0$. If $f(0)=1$ then
\begin{qoptions}
\opt{A}{Number of real roots of $f(x)-f(-x)=6x-10$ is 1}
\opt{B}{If $g(x)+f(x)=0$ then complete interval in which $g(x)$ is increasing is $x\in(-\infty,-1]\cup[1,\infty)$}
\opt{C}{The range of values of $K$ for which $f(x)=K$ has three distinct solutions is $K\in(0,2)$}
\opt{D}{Range of the function $P(\theta)=f(\sin\theta)$ where $\theta\in R$ is $[0,2]$}
\end{qoptions}
\end{bankq}
\begin{bankq}[id=TOAS08-P2-MATH-S1-Q06, source={TOAS - Test 08 (15 May 2022)}, paper=2, qno=6, type=multi, subtopic={Extrema}, answer={A,B,C}, answer_src=key, difficulty=7, diff_parts={K2 S4 I1 C0 F1}, diff_note={Quintic extrema give counting of real roots for a}]
Let $a\in R$ and $f:R\to R$ be given by $f(x)=3x^5-50x^3+135x-a$. The equation $f(x)=0$ will have :
\begin{qoptions}
\opt{A}{one real root for $|a|>216$}
\opt{B}{three real roots for $88<|a|<216$}
\opt{C}{five real roots for $|a|<88$}
\opt{D}{five real roots for $a\in R$}
\end{qoptions}
\end{bankq}
\begin{bankq}[id=TOAS08-P2-MATH-S2-Q07, source={TOAS - Test 08 (15 May 2022)}, paper=2, qno=7, type=single, subtopic={Increasing and decreasing functions}, answer={A}, answer_src=key, difficulty=8, diff_parts={K2 S4 I2 C1 F0}, diff_note={Convexity and monotonicity of phi; critical point conditions for g}]
Let $\phi(x)=3f\left(\dfrac{x^2}{3}\right)+f\left(3-x^2\right)\ \forall\,x\in(-3,4)$ where $f''(x)>0\ \forall\,x\in(-3,4)$,\\
$g(x)=(a^2-3a+2)\left(\cos^2\dfrac x4-\sin^2\dfrac x4\right)+(a-1)x+\sin1$\\
Match List-I with List-II and select the correct answer using the code given below the list.
\begin{center}\footnotesize\renewcommand{\arraystretch}{1.1}\setlength{\extrarowheight}{3.5pt}\setlength{\tabcolsep}{2pt}
\begin{tabular}{|c|l|c|l|}
\hline
\multicolumn{2}{|c|}{List-I} & \multicolumn{2}{c|}{List-II}\\[3pt]
\hline
(P) & $\phi(x)$ is increasing in & (1) & $(0,1)\cup(1,4)$\\[4pt]
(Q) & $\phi(x)$ is decreasing in & (2) & $\left[-\dfrac32,0\right]\cup\left[\dfrac32,4\right)$\\[6pt]
(R) & \parbox[t]{3.1cm}{The set of values of $a$ for which $g(x)$ does not possess critical point} & (3) & $\left(-3,-\dfrac32\right]\cup\left[0,\dfrac32\right]$\\[10pt]
(S) & \parbox[t]{3.1cm}{The set of values of $a$ for which $g(x)$ possess critical point} & (4) & $(-\infty,0]\cup\{1\}\cup[4,\infty)$\\[3pt]
\hline
\end{tabular}
\end{center}
\begin{qoptions}
\opt{A}{$P\to2;Q\to3;R\to1;S\to4$}
\opt{B}{$P\to2;Q\to3;R\to4;S\to1$}
\opt{C}{$P\to3;Q\to2;R\to1;S\to4$}
\opt{D}{$P\to3;Q\to2;R\to4;S\to1$}
\end{qoptions}
\end{bankq}
\begin{bankq}[id=TOAS09-P1-MATH-S1-Q02, source={TOAS - Test 09 (30 May 2022)}, paper=1, qno=2, type=single, subtopic={Extrema}, answer={D}, answer_src=key, difficulty=4, diff_parts={K2 S2 I1 C0 F0}, diff_note={Extrema: standard approach for this subtopic; 2 concept(s), 2 step(s), needs a non-obvious insight}]
The function $f(x)=\alpha x^2-\beta x-4x^3+\gamma$, where $\alpha,\beta,\gamma\in R$, has local minima at $P(\log_2a,\ f(\log_2a))$ \& local maxima at $Q(\log_2a^2,\ f(\log_2a^2))$. If the graph of $f(x)$ changes concavity about the point $R\!\left(\dfrac34,\ f\!\left(\dfrac34\right)\right)$, then which of the following conic section can have eccentricity `a' -
\begin{qoptions}
\opt{A}{circle}
\opt{B}{parabola}
\opt{C}{ellipse}
\opt{D}{hyperbola}
\end{qoptions}
\end{bankq}
\begin{bankq}[id=TOAS09-P1-MATH-S2-Q05, source={TOAS - Test 09 (30 May 2022)}, paper=1, qno=5, type=multi, subtopic={Extrema}, answer={A,C}, answer_src=key, difficulty=6, diff_parts={K2 S3 I1 C0 F1}, diff_note={Extrema: standard approach for this subtopic; 2 concept(s), 3 step(s), needs a non-obvious insight}]
Consider twice derivable bijective function defined on $[0,1]$ such that $f(\alpha)=\dfrac13$ and $f(f(\beta))=\dfrac13$. If $f''(x)>0\ \forall x\in\left[0,\dfrac12\right)$ and $f''(x)<0\ \forall x\in\left(\dfrac12,1\right]$ with $f\!\left(\dfrac12\right)=\dfrac12$. Which the following holds good ?
\begin{qoptions}
\opt{A}{$\beta>\alpha$ implies $f(x)$ is an increasing function}
\opt{B}{$\beta>\alpha$ implies $f(x)$ is a decreasing function}
\opt{C}{$2f\!\left(\dfrac14\right)<f(0)+f\!\left(\dfrac12\right)$}
\opt{D}{$2f\!\left(\dfrac34\right)<f\!\left(\dfrac12\right)+f(1)$}
\end{qoptions}
\end{bankq}
\begin{bankq}[id=TOAS09-P1-MATH-S2-Q06, source={TOAS - Test 09 (30 May 2022)}, paper=1, qno=6, type=multi, subtopic={Extrema}, answer={A,B}, answer_src=key, difficulty=6, diff_parts={K2 S3 I1 C0 F1}, diff_note={Extrema: standard approach for this subtopic; 2 concept(s), 3 step(s), needs a non-obvious insight}]
For the function $f(x)=e^{e^x}-\tan^{-1}x-3$ which of the following holds good ?
\begin{qoptions}
\opt{A}{$f(x)=0$ has exactly one positive real root}
\opt{B}{$y=f(x)$ has exactly one minima at negative $x$}
\opt{C}{$y=f(x)$ has exactly one point of minima at positive $x$}
\opt{D}{$y=f(x)$ is strictly increasing function}
\end{qoptions}
\end{bankq}
\begin{bankq}[id=TOAS09-P1-MATH-S2-Q08, source={TOAS - Test 09 (30 May 2022)}, paper=1, qno=8, type=multi, subtopic={Increasing and decreasing functions}, answer={B,C}, answer_src=key, difficulty=2, diff_parts={K1 S1 I0 C0 F1}, diff_note={Increasing and decreasing functions: standard approach for this subtopic; 1 concept(s), 1 step(s)}]
Let $f(x)$ be a non-negative continuous and bounded function for all $x\geqslant0$. If $(\cos x)f'(x)\leqslant(\sin x-\cos x)f(x)$ ; $\forall x\geqslant0$, then which of the following is / are correct ?
\begin{qoptions}
\opt{A}{$f(6)+f(5)>0$}
\opt{B}{$x^2-3x+2+f(7)=0$ has 2 distinct solutions}
\opt{C}{$f(5)f(7)-f(6)f(5)=0$}
\opt{D}{$\displaystyle\lim_{x\to6}\dfrac{f(x)-\sin(\pi x)}{(x-6)}=1$}
\end{qoptions}
\end{bankq}
\begin{bankq}[id=TOAS09-P2-MATH-S1-Q06, source={TOAS - Test 09 (30 May 2022)}, paper=2, qno=6, type=single, subtopic={Extrema}, answer={D}, answer_src=key, difficulty=1, diff_parts={K1 S1 I0 C0 F0}, diff_note={Extrema: standard approach for this subtopic; 1 concept(s), 1 step(s)}]
The number of real solutions of the equation $\tan(e^x)=e^x+e^{-x}$, $x>0$, is
\begin{qoptions}
\opt{A}{0}
\opt{B}{1}
\opt{C}{2}
\opt{D}{none of these}
\end{qoptions}
\end{bankq}
\begin{bankq}[id=TOAS09-P2-MATH-S1-Q10, source={TOAS - Test 09 (30 May 2022)}, paper=2, qno=10, type=single, subtopic={Extrema}, answer={C}, answer_src=key, difficulty=2, diff_parts={K1 S2 I0 C0 F0}, diff_note={Extrema: standard approach for this subtopic; 1 concept(s), 2 step(s)}]
Let $f:R\to R$, $f(x)=ax^3+bx^2+cx+d$, has no extreme value, then which of the following is always correct ?
\begin{qoptions}
\opt{A}{$3a+2b+c\geq0$}
\opt{B}{$c\leq0$}
\opt{C}{$(3a+2b+c)c\geq0$}
\opt{D}{$3a+2b+2c\geq0$}
\end{qoptions}
\end{bankq}
\begin{bankq}[id=TOAS09-P2-MATH-S1-Q16, source={TOAS - Test 09 (30 May 2022)}, paper=2, qno=16, type=single, subtopic={Extrema}, answer={D}, answer_src=key, difficulty=4, diff_parts={K2 S2 I1 C0 F0}, diff_note={Extrema: standard approach for this subtopic; 2 concept(s), 2 step(s), needs a non-obvious insight}]
In the interval $(0,1)$ maximum value of the function $f(x)=|x\,\ell nx|$ is
\begin{qoptions}
\opt{A}{0}
\opt{B}{1}
\opt{C}{$e$}
\opt{D}{$e^{-1}$}
\end{qoptions}
\end{bankq}
\begin{bankq}[id=TOAS09-P2-MATH-S2-Q01, source={TOAS - Test 09 (30 May 2022)}, paper=2, qno=1, type=numerical, subtopic={Optimisation}, answer={6}, answer_src=key, difficulty=8, diff_parts={K3 S3 I1 C1 F1}, diff_note={Optimisation: standard approach for this subtopic; 3 concept(s), 3 step(s), needs a non-obvious insight, heavier computation}]
If $m$ is the least value of $f(x)=\sqrt{x^2-2x+2}+\sqrt{x^2-4x+29}$, then $[m]$ is equal to\par
(where $[.]$ denotes greatest integer function)
\end{bankq}
\begin{bankq}[id=TOAS09-P2-MATH-S2-Q02, source={TOAS - Test 09 (30 May 2022)}, paper=2, qno=2, type=numerical, subtopic={Extrema}, answer={5}, answer_src=key, difficulty=7, diff_parts={K2 S3 I1 C1 F1}, diff_note={Extrema: standard approach for this subtopic; 2 concept(s), 3 step(s), needs a non-obvious insight, heavier computation}]
If `a' and `b' are the greatest and least values of $f(x)=18x^{4/3}-9x^{1/3}+1$ in the interval $[-1,8]$ then the value of $(a+112b)$ is
\end{bankq}
\begin{bankq}[id=TOAS09-P2-MATH-S2-Q06, source={TOAS - Test 09 (30 May 2022)}, paper=2, qno=6, type=numerical, subtopic={Extrema}, answer={5}, answer_src=solved, working={Maxima at 0 and 1 need $-1\le a<1$: integers $-1,0$ so $L=2$; extrema at $-1$ (min), 0, 1: $M=3$; sum 5}, difficulty=7, diff_parts={K2 S3 I1 C1 F1}, diff_note={Extrema: standard approach for this subtopic; 2 concept(s), 3 step(s), needs a non-obvious insight, heavier computation}]
$f(x)=\begin{cases}(x+1)^2 & -3<x\leqslant0\\ 2x^{1/3}+a & 0<x\leqslant1\\ 2-x & 1<x<2\end{cases}$ has maxima at $x=0,1$. If number of integral values of `a' is $L$ and total number of extrema of $f(x)$ is $M$ then $(L+M)$ is
\end{bankq}
\end{document}
